\chapter{Quantum Machine Learning}
\label{ch:qml}

This chapter develops the quantum and hybrid model families that
Chapters~\ref{ch:methodology} and \ref{ch:experiments} put to the test. The
running theme is that every design choice --- how data enter the circuit, how
the trainable part is shaped, what is measured, and how parameters are
optimized --- has consequences that are now reasonably well understood
theoretically, and that these consequences (expressivity spectra, kernel
equivalences, barren plateaus, concentration) directly generate the hypotheses
of the experimental programme.

\section{A taxonomy of quantum machine learning}
\label{sec:qml-taxonomy}

Following the scheme popularized by \citet{schuld2021book}, QML approaches are
classified by whether the \emph{data} and the \emph{processing device} are
classical (C) or quantum (Q). \textbf{CC} is classical ML, possibly
quantum-inspired. \textbf{CQ} --- classical data, quantum processing --- is
the subject of this thesis: feature vectors are encoded into quantum states,
processed by circuits, and read out by measurement. \textbf{QC} uses
classical ML to support quantum experiments (calibration, tomography,
control), and \textbf{QQ} learns on inherently quantum data, the setting where
the strongest advantage arguments live \citep{huang2022advantage}; both are
outside the experimental scope and appear only in the outlook.

Within CQ, this thesis studies the two dominant paradigms --- \emph{quantum
kernel methods} (Section~\ref{sec:qml-kernels}) and \emph{variational quantum
circuits} (Sections~\ref{sec:qml-vqa}--\ref{sec:qml-architectures}) --- plus
\emph{hybrid architectures} that compose them with classical networks
(Section~\ref{sec:qml-hybrid}). The taxonomy earns its place beyond
bookkeeping with one observation: a quantum kernel inside a classical SVM is
already a hybrid CQ pipeline, so ``quantum versus hybrid versus classical'' is
a spectrum of composition choices, not three disjoint species.

\section{Data encoding}
\label{sec:qml-encoding}

A CQ model must first map $\bm{x} \in \R^d$ to an $n$-qubit state
$\ket{\phi(\bm{x})} = \enc{\bm{x}} \ket{0}^{\otimes n}$ via an encoding
(feature-map) circuit $\enc{\bm{x}}$. This choice is not preprocessing
trivia: it fixes the geometry of the data in Hilbert space, hence the kernel
the model can realize and the function class it can express
\citep{schuld2019hilbert, schuld2021fourier, larose2020encodings}.

\paragraph{Basis encoding} writes a binary string into the computational
basis, $\bm{b} \mapsto \ket{\bm{b}}$. It is qubit-hungry for continuous data
and is used here only conceptually.

\paragraph{Angle encoding} uses one feature per qubit as a rotation angle,
\begin{equation}
  \enc{\bm{x}} = \bigotimes_{j=1}^{d} R_{Y}(x_j),
  \qquad
  \ket{\phi(\bm{x})} = \bigotimes_{j=1}^{d}
    \Bigl( \cos\tfrac{x_j}{2}\ket{0} + \sin\tfrac{x_j}{2}\ket{1} \Bigr),
  \label{eq:angle-encoding}
\end{equation}
requiring $n = d$ qubits and depth 1. Because the rotation is $2\pi$-periodic,
features must be confined to a fixed interval to avoid wrap-around aliasing;
this motivates the thesis-wide convention of min--max scaling every feature to
$[0, \pi]$, fit on training folds only (Chapter~\ref{ch:methodology}).
Unentangled angle encoding produces a product state, and its induced kernel
factorizes into a product of per-feature $\cos$ terms --- a useful reminder
that ``quantum'' does not yet mean ``hard to simulate.''

\paragraph{Entangled feature maps.} Expressive power beyond products comes
from interleaving encodings with entanglers. The canonical example is the
ZZ (IQP-style) feature map of \citet{havlicek2019kernels},
\begin{equation}
  \enc{\bm{x}}
  = \Bigl[ U_{\Phi}(\bm{x})\, H^{\otimes n} \Bigr]^{L},
  \qquad
  U_{\Phi}(\bm{x}) = \exp\Bigl( i \sum_{j} \phi_j(\bm{x})\, Z_j
      + i \sum_{j < k} \phi_{jk}(\bm{x})\, Z_j Z_k \Bigr),
  \label{eq:zz-map}
\end{equation}
with the standard choices $\phi_j(\bm{x}) = x_j$ and
$\phi_{jk}(\bm{x}) = (\pi - x_j)(\pi - x_k)$, repeated $L$ times. The
product terms couple pairs of features through entangling phases ---
conjectured hard to simulate classically at depth $L \ge 2$ --- and make
parity-like feature interactions natively available, which is precisely why
the sign-parity benchmark task of Chapter~\ref{ch:methodology} is the
designated stress test for this family.

\paragraph{Amplitude encoding} stores $\bm{x}$ (padded, normalized) in the
amplitudes of \eqref{eq:nqubit}: $2^n$ features fit in $n$ qubits, but state
preparation generally costs depth exponential in $n$ and the encoding is
nonlinear in $\bm{x}$ through normalization. It appears in this thesis only
as a comparison point in the encoding ablation (E4).

\paragraph{Data re-uploading} repeats the encoding several times, interleaved
with trainable blocks \citep{perezsalinas2020reuploading}:
\begin{equation}
  \ket{\phi_{\bm{\theta}}(\bm{x})}
  = \prod_{l=1}^{L} \Bigl[ W_l(\bm{\theta}_l)\, \enc{\bm{x}} \Bigr]
    \ket{0}^{\otimes n}.
  \label{eq:reuploading}
\end{equation}
Re-uploading is not a redundancy: it provably enlarges the model's frequency
content, as the following Fourier picture makes precise.

\paragraph{The Fourier lens on expressivity.}
\citet{schuld2021fourier} showed that any model of the form
$f_{\bm{\theta}}(\bm{x}) = \bra{\phi_{\bm{\theta}}(\bm{x})} M
\ket{\phi_{\bm{\theta}}(\bm{x})}$ whose encoding gates are generated by fixed
Hamiltonians, $e^{-i x_j H_j}$, is a generalized trigonometric polynomial
\begin{equation}
  f_{\bm{\theta}}(\bm{x})
  = \sum_{\bm{\omega} \in \Omega} c_{\bm{\omega}}(\bm{\theta}, M)\,
    e^{i \bm{\omega}^{\top} \bm{x}},
  \label{eq:fourier}
\end{equation}
whose frequency set $\Omega$ is determined \emph{entirely by the encoding} ---
the pairwise differences of eigenvalue sums of the $H_j$ --- while the
trainable parts and the observable only shape the coefficients
$c_{\bm{\omega}}$. For Pauli-rotation encodings, one upload per feature yields
frequencies $\{-1, 0, 1\}$ per dimension; $L$ uploads extend this to
$\{-L, \dots, L\}$. Three consequences organize much of what follows:
(i) a single-upload angle-encoded model is, in each feature, a degree-one
trigonometric function --- expressivity claims must be checked against this
sobering baseline; (ii) re-uploading, not deeper ansätze, is the lever that
adds frequencies; and (iii) expressivity is decoupled from trainability ---
enlarging $\Omega$ says nothing about whether $c_{\bm{\omega}}$ can be
efficiently fit (Section~\ref{sec:qml-trainability}). Experiment E4
demonstrates \eqref{eq:fourier} directly by fitting one-dimensional target
functions with increasing upload counts.

\section{Variational quantum algorithms}
\label{sec:qml-vqa}

The variational paradigm \citep{cerezo2021vqa} builds models from three
components: an encoding $\enc{\bm{x}}$, a trainable ansatz $\ansatz$, and a
measured observable $M$,
\begin{equation}
  f_{\bm{\theta}}(\bm{x})
  = \bra{0^{\otimes n}} \enc{\bm{x}}^{\dagger}\, W^{\dagger}(\bm{\theta})\,
    M\, W(\bm{\theta})\, \enc{\bm{x}} \ket{0^{\otimes n}},
  \label{eq:vqc-model}
\end{equation}
trained by a classical optimizer minimizing an empirical loss
$C(\bm{\theta}) = \tfrac{1}{m} \sum_i \ell\bigl( f_{\bm{\theta}}(\bm{x}_i),
y_i \bigr)$ --- a quantum--classical loop in which the quantum device only
evaluates \eqref{eq:vqc-model}. In this thesis $M$ is a single-qubit
$Z$ observable unless stated otherwise, making
$f_{\bm{\theta}} \in [-1, 1]$ a signed classification score
\citep{mitarai2018qcl, farhi2018qnn, schuld2020circuitcentric}.

\paragraph{Ansatz families.} The experiments use two representative ansätze:
\emph{hardware-efficient} layers of single-qubit rotations plus a fixed
entangling pattern \citep{kandala2017hardware}, chosen for shallow depth on
real devices; and \emph{strongly entangling layers}
\citep{schuld2020circuitcentric}, with all-to-all-leaning connectivity and
three rotation angles per qubit per layer. Depth is swept explicitly (E3),
since it controls both capacity and every pathology of
Section~\ref{sec:qml-trainability}.

\paragraph{Exact gradients: the parameter-shift rule.}
For gates of the form \eqref{eq:rotation} --- generators with two eigenvalues
$\pm\tfrac{1}{2}$ --- the derivative of \eqref{eq:vqc-model} with respect to
the $k$-th angle is \emph{exactly}
\begin{equation}
  \frac{\partial f_{\bm{\theta}}}{\partial \theta_k}
  = \frac{1}{2} \Bigl[
      f_{\bm{\theta} + \frac{\pi}{2} \bm{e}_k}
    - f_{\bm{\theta} - \frac{\pi}{2} \bm{e}_k}
    \Bigr],
  \label{eq:parameter-shift}
\end{equation}
two evaluations of the same circuit at shifted parameters
\citep{mitarai2018qcl, schuld2021fourier} --- not a finite-difference
approximation, and compatible with hardware since only expectation values are
needed. The cost, however, scales linearly in the parameter count \emph{per
gradient step} (two circuit evaluations per parameter, each itself averaged
over shots), which is why gradient-based VQC training on hardware is
prohibitive and why this thesis trains exclusively in simulation, reserving
hardware for inference (Chapter~\ref{ch:methodology}).

\paragraph{Optimizers.} Experiment E3 compares Adam \citep{kingma2015adam} on
parameter-shift gradients with two gradient-light alternatives standard in
the VQA literature: SPSA \citep{spall1992spsa}, which estimates a descent
direction from two evaluations regardless of dimension and is naturally robust
to shot noise, and the derivative-free trust-region method COBYLA
\citep{powell1994cobyla}. The optimizer axis is treated as a first-class
design factor (RQ2), not a nuisance setting.

\section{Quantum kernel methods}
\label{sec:qml-kernels}

The encoding circuit alone already defines a learning method. The
\emph{fidelity quantum kernel} associated with $\enc{\cdot}$ is
\begin{equation}
  \qkernel(\bm{x}, \bm{x}')
  = \abs{ \braket{\phi(\bm{x})}{\phi(\bm{x}')} }^2
  = \abs{ \bra{0^{\otimes n}} \enc{\bm{x}}^{\dagger} \enc{\bm{x}'}
          \ket{0^{\otimes n}} }^2,
  \label{eq:fidelity-kernel}
\end{equation}
a positive-semidefinite similarity measure computable on a quantum device by
preparing $\enc{\bm{x}}^{\dagger} \enc{\bm{x}'} \ket{0^{\otimes n}}$ and
measuring the probability of returning to $\ket{0^{\otimes n}}$ (the
``inversion test''), at circuit depth twice that of the encoding
\citep{havlicek2019kernels, schuld2019hilbert}. The QSVM pipeline estimates
the Gram matrix $K_{ij} = \qkernel(\bm{x}_i, \bm{x}_j)$ --- $O(m^2)$ circuit
families, each averaged over shots --- and hands it to the convex classical
SVM dual \eqref{eq:svm-dual}. Against the RBF-SVM baseline
\eqref{eq:rbf}, the QSVM differs in exactly one component, which is what
makes the E2 comparison clean.

\paragraph{Quantum models are kernel methods.}
\citet{schuld2021kernels} sharpened the connection: the class of models
\eqref{eq:vqc-model}, viewed as functions of the encoded density matrix
$\rho(\bm{x}) = \ketbra{\phi(\bm{x})}{\phi(\bm{x})}$, is \emph{linear} in
$\rho(\bm{x})$, i.e., a linear model in the feature space of density matrices
with inner product $\Tr[\rho(\bm{x}) \rho(\bm{x}')] = \qkernel(\bm{x},
\bm{x}')$. By representer-theorem reasoning, the loss-optimal model in that
class is a kernel expansion over the training data --- so a trained VQC can,
at best, match the regularized kernel machine built from
\eqref{eq:fidelity-kernel}, and variational training is better understood as
implicit (possibly advantageous, possibly lossy) regularization within the
kernel's hypothesis space. This result frames E2 and E3 as two views of one
object and cautions against narratives in which the ansatz ``adds power''
beyond the encoding.

\paragraph{Kernel diagnostics.} Two quantities computed in E2 connect kernels
to outcomes before any SVM is trained: \emph{kernel--target alignment}, the
normalized inner product between the Gram matrix and the ideal label kernel
$\bm{\tilde{y}} \bm{\tilde{y}}^{\top}$, a cheap predictor of learnability; and
the \emph{spectrum} of $K$, whose flatness signals the concentration pathology
below. The \emph{projected quantum kernel} of \citet{huang2021power} ---
built from reduced single-qubit density matrices,
$\kappa^{\mathrm{PQ}}(\bm{x}, \bm{x}') = \exp\bigl( -\gamma \sum_k
\norm{\rho_k(\bm{x}) - \rho_k(\bm{x}')}_F^2 \bigr)$ --- is included as the
literature's main remedy: by projecting before comparing, it tempers the
exponential-dimension geometry that drives both concentration and poor
generalization \citep{kubler2021inductive}.

\section{Model architectures}
\label{sec:qml-architectures}

Three circuit architectures span the variational side of the comparison.

\paragraph{VQC.} The plain variational quantum classifier
\eqref{eq:vqc-model}, instantiated over the grid \{angle, ZZ\} encodings
$\times$ \{hardware-efficient, strongly entangling\} ansätze $\times$ depth
$\{1, 3, 5\}$ --- the E3 design. Parameter counts ($3nd_{\mathrm{ansatz}}$ for
strongly entangling layers on $n$ qubits) are logged for every run and define
the MLP matching budgets.

\paragraph{Data re-uploading classifier.} The single- or few-qubit model
\eqref{eq:reuploading} of \citet{perezsalinas2020reuploading}, universal for
classification on one qubit in the many-upload limit and, per
\eqref{eq:fourier}, the cleanest experimental handle on the
frequency-enrichment story (E4). It is also among the strongest small quantum
models in independent benchmarks \citep{bowles2024better}, at minimal
simulation cost.

\paragraph{QCNN.} The quantum convolutional network of \citet{cong2019qcnn}
alternates translationally-shared two-qubit ``convolution'' unitaries with
measurement-free ``pooling'' that discards half the qubits per layer, giving
$O(\log n)$ depth and parameter counts logarithmic in the input size; the
classical-data variant follows \citet{hur2022qcnn}. QCNNs are the image-tier
quantum contender (8 qubits over PCA-compressed pixel features) and are of
special theoretical interest because the architecture is among the few with
evidence of \emph{avoiding} barren plateaus \citep{larocca2025barren} --- at
the price of falling within classically simulable families
\citep{cerezo2023simulability}, a tension E5/E9 quantify in practice.

\section{Trainability: barren plateaus and concentration}
\label{sec:qml-trainability}

The central obstruction to scaling variational models is not expressivity but
\emph{trainability}. \citet{mcclean2018barren} showed that for ansätze random
enough to approximate unitary 2-designs, gradients of cost functions like
\eqref{eq:vqc-model} vanish in probability:
\begin{equation}
  \E\bigl[ \partial_{\theta_k} C \bigr] = 0,
  \qquad
  \Var\bigl[ \partial_{\theta_k} C \bigr] \in O\!\bigl( b^{-n} \bigr)
  \text{ for some } b > 1,
  \label{eq:barren}
\end{equation}
the \emph{barren plateau}: with overwhelming probability over random
initializations, the loss landscape is exponentially flat in the qubit count,
and estimating a useful gradient direction through shot noise
\eqref{eq:shotnoise} costs exponentially many measurements.
\citet{cerezo2021costfunction} refined the picture: with shallow
(logarithmic-depth) ansätze, \emph{local} observables (as used throughout this
thesis) admit gradients decaying at worst polynomially, while \emph{global}
observables plateau at any depth --- making cost design a first-order modeling
decision. The modern synthesis \citep{larocca2025barren} identifies
randomness, entanglement growth, noise, and cost globality as interchangeable
sources of the same exponential concentration; strategies that provably escape
it (structured ansätze such as QCNNs, clever initializations, problem-tailored
symmetry) tend to land in circuit families that are classically simulable
\citep{cerezo2023simulability} --- arguably the sharpest form of the
quantum-utility dilemma. Experiment E6 reproduces \eqref{eq:barren} directly,
measuring gradient variance versus qubit count ($2$--$16$) and depth, for
local versus global costs.

Quantum kernels suffer a mirror-image pathology: \emph{exponential
concentration} of kernel values \citep{thanasilp2024concentration}. For
sufficiently expressive (or deep, or noisy) encodings, $\qkernel(\bm{x},
\bm{x}')$ concentrates around a data-independent constant with deviations
exponentially small in $n$ --- so the Gram matrix approaches a scaled
identity, generalization collapses, and resolving the residual signal requires
exponentially many shots. Bandwidth-style fixes (rescaling feature ranges) and
projected kernels trade expressivity back for signal
\citep{kubler2021inductive, huang2021power}. E2 measures kernel-value spread
versus qubit count on this thesis's datasets. Finally, on the generalization
side, uniform bounds scaling with trainable-gate counts \citep{caro2022generalization}
--- and their empirical refutation as a complete story
\citep{gilfuster2024rethinking} --- plus the effective-dimension perspective
\citep{abbas2021power} justify treating capacity claims empirically (RQ5)
rather than as settled theory.

\section{Hybrid architectures}
\label{sec:qml-hybrid}

Hybrid models compose classical and quantum components under end-to-end or
staged training; automatic differentiation through quantum nodes
\citep{bergholm2018pennylane} makes the composition practical. Three designs
enter the comparison (E5), each with a matched classical control.

\paragraph{Compression front-ends.} PCA (or a small autoencoder) reduces
$d$-dimensional inputs to the qubit budget before a VQC. This is the standard
route for image data and doubles as the fair way to give quantum models
high-dimensional inputs at all; the control is the same compressed features
fed to classical heads.

\paragraph{Classical feature extractor with quantum head.} A small CNN
produces $k$ features consumed by a $k$-qubit VQC head, trained jointly; the
control replaces the quantum head with a dense head parameter-matched to the
VQC's angle count. This is the cleanest hybrid-versus-classical contrast in
the thesis: identical data, identical backbone, matched capacity, single
varying component.

\paragraph{Quantum transfer learning.} Following \citet{mari2020transfer}, a
frozen pretrained backbone feeds a trainable ``dressed quantum circuit'' ---
classical linear map in, PQC, classical linear map out. Notably, this family
ranked among the stronger quantum-adjacent models in independent benchmarking
\citep{bowles2024better}, with the open question --- how much of its
performance is the dressing rather than the quantum core --- addressed by the
parameter-matched dressed-\emph{classical} control.

Throughout, the taxonomy of Section~\ref{sec:qml-taxonomy} keeps the
bookkeeping honest: the quantum kernel SVM is a CQ hybrid with quantum
similarity and classical optimization; the dressed circuit is a CQ hybrid with
classical representation learning and quantum decision surface. What the
experiments compare is therefore not ``quantum versus classical'' as slogans,
but specific placements of a quantum component inside otherwise-matched
pipelines --- exactly the resolution at which RQ3 is answerable.
